%!TEX root = ../dolgozat.tex
\pagestyle{empty}

\section*{Adathordozó használati útmutató}

A szakdolgozathoz mellékelt adathordozó a következő könyvtárakat és fájlokat tartalmazza:
\begin{itemize}
\item \texttt{dolgozat.pdf} -- a szakdolgozat PDF formátumban,
\item \texttt{dolgozat/} -- a szakdolgozat \LaTeX{} forráskódja,
\item \texttt{online-exam-system/backend/} -- a FastAPI alapú szerveroldali alkalmazás,
\item \texttt{online-exam-system/frontend/} -- a Next.js alapú kliensoldali alkalmazás,
\item \texttt{online-exam-system/docker-compose.yml} -- a PostgreSQL adatbázis indítását leíró fájl,
\item \texttt{online-exam-system/README.md} -- részletes telepítési és használati leírás,
\item \texttt{tesztek/} -- a \mref{ch:teszteles}ben ismertetett tesztszkriptek.
\end{itemize}

\subsection*{Futtatás}

A futtatáshoz Docker Desktop, Python 3.10 vagy újabb, valamint Node.js 18 vagy újabb verzió
szükséges. Az alkalmazás három lépésben indítható el, az \texttt{online-exam-system} könyvtárból:
\begin{enumerate}
\item Az adatbázis és a sandbox image előkészítése:\
  \texttt{docker compose up -d}\
  \texttt{docker pull python:3.11-slim}
\item A szerver indítása a \texttt{backend} könyvtárban:\
  \texttt{python -m venv venv}, majd a környezet aktiválása után\
  \texttt{pip install -r requirements.txt}\
  \texttt{uvicorn main:app --port 8000}
\item A kliens indítása a \texttt{frontend} könyvtárban:\
  \texttt{npm install}\
  \texttt{npm run dev}
\end{enumerate}
Ezután az alkalmazás a \texttt{http://localhost:3000} címen érhető el, az API interaktív
dokumentációja pedig a \texttt{http://localhost:8000/docs} címen. Oktatói fiók a regisztrációs
űrlapon az ,,Oktatóként regisztrálok'' lehetőséggel és az alapértelmezett \texttt{oktato123}
meghívókóddal hozható létre (ez a \texttt{TEACHER\_INVITE\_CODE} környezeti változóval
módosítható).

A tesztszkriptek a \texttt{backend} könyvtár virtuális környezetének Python értelmezőjével
futtathatók; paraméterként a \texttt{backend} könyvtár elérési útját kell megadni. A sandbox
tesztjéhez futó Docker szükséges.
